\documentclass{article}
\usepackage{amsmath,amssymb}
\usepackage{xurl}
\usepackage[hidelinks]{hyperref}
% Shared commands for notes in docs/paper-gaps/.

\DeclareUnicodeCharacter{2080}{\ifmmode{}_{0}\else\textsubscript{0}\fi}
\DeclareUnicodeCharacter{2081}{\ifmmode{}_{1}\else\textsubscript{1}\fi}
\DeclareUnicodeCharacter{2082}{\ifmmode{}_{2}\else\textsubscript{2}\fi}
\DeclareUnicodeCharacter{2099}{\ifmmode{}_{n}\else\textsubscript{n}\fi}

\newcommand{\Lean}{\textsc{Lean}}
\ExplSyntaxOn
\NewDocumentCommand{\leanid}{m}{
  \ifmmode
    \textsf{\detokenize{#1}}
  \else
    \group_begin:
    \urlstyle{sf}
    \tl_set:Nn \l_tmpa_tl {#1}
    \tl_replace_all:Nnn \l_tmpa_tl {\_} {_}
    \exp_args:NV \path \l_tmpa_tl
    \group_end:
  \fi
}
\ExplSyntaxOff
\providecommand{\tr}{\operatorname{tr}}
\newcommand{\tnleanIssueBase}{https://github.com/LionSR/TNLean/issues/}
\newcommand{\tnleanPrBase}{https://github.com/LionSR/TNLean/pull/}
\newcommand{\tnleanDocBase}{https://sirui-lu.com/TNLean/docs/}
\newcommand{\ghissue}[1]{\href{\tnleanIssueBase#1}{GitHub issue \##1}}
\newcommand{\ghpr}[1]{\href{\tnleanPrBase#1}{PR \##1}}
\newcommand{\doclink}[1]{\href{\tnleanDocBase#1.html}{\path{#1}}}

% Dirac notation and the identity operator, as in blueprint/src/macros/common.tex.
% \providecommand keeps note-local definitions authoritative.
\usepackage{dsfont}
\providecommand{\ket}[1]{|#1\rangle}
\providecommand{\bra}[1]{\langle #1|}
\providecommand{\braket}[2]{\langle #1 | #2 \rangle}
\providecommand{\ketbra}[2]{|#1\rangle\!\langle #2|}
\providecommand{\Id}{\mathds{1}}
\providecommand{\one}{\mathds{1}}
\providecommand{\C}{\mathbb{C}}
\providecommand{\MN}[1]{M_{#1}(\C)}
\providecommand{\MD}{M_D(\C)}

% Verdict marker: \gapnote{<kind>}{<status>}. Kind is the policy
% classification (clarification, local-correction, scope-restriction,
% unfaithful, false-source, open-gap); status is open, wip (elimination
% underway), resolved, or historical. Machine-read by the site index;
% prints a verdict line.
\newcommand{\gapnote}[2]{\par\noindent\textbf{Verdict:} #1 (#2).\par}

\title{The Tree Circuit with Measurements: Rounds, Registers, and Chain Length}
\author{}
\date{2026-10-01}
\begin{document}
\maketitle
\gapnote{scope-restriction}{open}

\section{What the source states}
The paragraph ``Tree-RG circuit with measurements''
of~\cite{Malz2024Preparation} observes that the isometries of the tree circuit,
eq.~(16), ``act on a constant number of sites which, although spatially
separated, can be teleported at neighboring registers with a constant
overhead'', by ``creating nearest-neighbor entangled pairs, then performing
simultaneous measurements, and correcting (without postselection) based on the
measurement outcomes''. Hence ``every isometry in eq.~(16) takes constant time
using measurement'', the tree circuit has $O(\log\log(N/\varepsilon))$ layers,
and, with the fixed-point state prepared in constant depth, a short-range
correlated matrix product state is prepared with measurements in depth
$O(\log\log(N/\varepsilon))$. The source fixes no model of circuits with
measurements.

\section{What is formalized}
\begin{enumerate}
    \item \emph{The model.} A measurement round is a local circuit, a
          measurement of some sites in the computational basis, and a product of
          single-site unitaries chosen from the outcomes; a preparation is a
          sequence of rounds, the depth being the total number of layers of
          their circuits. Every round is a protocol of
          $\mathrm{QCcc}_\ell$ of~\cite{Piroli2021LOCC} of no larger depth;
          that paper performs a single round and remarks that ``one could also
          define a more general scheme with multiple rounds of LOCC''. A
          sequence of rounds is a composition of such transformations. The
          number of rounds is not bounded: the tree with $k$ levels uses $2k$
          rounds, whereas in the class $\mathrm{QCcc}^{(k)}_\ell$ of the
          paragraph ``Phases of matter'' of~\cite{Piroli2021LOCC} the number
          $k$ of composed transformations does not depend on the system
          size. Several rounds are needed because a
          teleported register must be corrected before the next isometry acts on
          it.
    \item \emph{Teleportation.} For qudits of any dimension $d$, any number of
          registers is moved along chains of nearest-neighbour hops of any
          length, in one round of depth $2$; for every outcome the round acts
          as $d^{-H}$ times the permutation of sites, $H$ the number of hops.
    \item \emph{One layer.} Two-site gates on pairwise disjoint stretches
          $a,\dots,a+2L+1$, with the sites strictly inside each stretch in
          $\ket0$, are applied in two rounds of total depth $5$, whatever $L$.
          A gate joins the sites $a$ and $a+2L+1$, at odd separation along the
          ring; gates at even separation, or on overlapping stretches, are not
          covered, and the tree needs neither.
    \item \emph{The tree.} A binary tree of two-site gates with $k$ levels, on the
          complete blocks of $2^k$ sites and on the complete smaller blocks
          after them, is applied with measurements in depth $5k$, on the vectors
          with $\ket0$ strictly inside the complete blocks of every level.
    \item \emph{Registers of several sites.} A register is $s$ consecutive
          sites. A unitary on two registers $a,\dots,a+s-1$ and
          $a+2L+s,\dots,a+2L+2s-1$ is applied by teleporting the first register
          site by site, the site $a+s-1$ first, in $s$ rounds of depth $2$,
          applying the unitary to $2s$ consecutive sites by a local circuit of
          depth $K$, and teleporting back; a layer of such unitaries on disjoint
          stretches takes depth $4s+K+2$, whatever $L$. Trees of such unitaries
          on the $M$ blocks of $q=s2^{k+1}$ sites take depth $(k+1)(4s+K+2)$.
    \item \emph{The approximating state.} If the blocked tensor of $A$ over $s$
          sites is injective, a register carries $\mathbb{C}^{D^2}$, every layer
          of the tree factorization of the blocked tensor over $s$ sites is an
          isometry, and the tree of unitaries extending these isometries
          implements the isometry $V$ of the polar decomposition of $A$ blocked
          over $q$ sites. After the pairs of the fixed point, prepared by a local
          circuit, the protocol outputs the approximating state of eq.~(10) with
          blocks of $q$ sites, in depth at most $C(k+1)$ with $C$ depending on $d$
          and $s$.
    \item \emph{The depth.} For a normal tensor, $s$ is a length at which the
          blocked tensor in the gauge of eq.~(5) is injective, and with
          Lemma~1'(i) at slope $a=\xi$ the error is at most $\varepsilon$ once
          $s2^{k+1}\ge a\log(N/\varepsilon)+b$. If moreover
          $s2^{k+1}\le2(a\log(N/\varepsilon)+b)$, then
          $k\le\log_2(a\log(N/\varepsilon)+b)$ and the depth is
          $O(\log\log(N/\varepsilon))$.
\end{enumerate}
The blocking over $s$ sites makes the two-site injectivity assumed for the
multiscale network of~\cite{gap:mswc24_tree_mera_scope} unnecessary here: the
tree is that of the blocked tensor, whose sites are the registers.
The formal statements\footnote{Formal declarations:
    \leanid{MPSPreparation.TeleportHop.isImplementationOn_round} in
    \doclink{TNLean/MPS/Preparation/TeleportationRound},
    \leanid{MPSPreparation.LongRangeGate.isRoundsImplementationOn_rounds} in
    \doclink{TNLean/MPS/Preparation/LongRangeGates}, and
    \leanid{MPSPreparation.isPreparedWithMeasurementRoundsInDepth_treeOp} in
    \doclink{TNLean/MPS/Preparation/TreeMeasurement},
    \leanid{MPSPreparation.RegisterGate.exists_rounds} in
    \doclink{TNLean/MPS/Preparation/RegisterGates},
    \leanid{MPSPreparation.exists_rounds_blockLayerOp_regTreeOp} in
    \doclink{TNLean/MPS/Preparation/RegisterTree},
    \leanid{MPSPreparation.regTreeOp_apply_blockInputCfg} in
    \doclink{TNLean/MPS/Preparation/RegisterTreeState}, and
    \leanid{MPSPreparation.exists_isPreparedWithMeasurementRoundsInDepth_approximatingMPVState},
    \leanid{MPSPreparation.exists_isPreparedWithMeasurementRoundsInDepth_approximationError_le}
    and
    \leanid{MPSPreparation.exists_isPreparedWithMeasurementRoundsInDepth_le_logb_log}
    in \doclink{TNLean/MPS/Preparation/LogLogDepthPreparation}.}
are exact, and no outcome is post-selected.

\section{The restriction}
The block length $q=s2^{k+1}$ divides the chain length $N$, so that the $N/q$
blocks of eq.~(10) all have $q$ sites and every tree is complete. The depth
bound therefore holds for the chain lengths $N$ divisible by $s2^{k+1}$ for some
$k$ with $a\log(N/\varepsilon)+b\le s2^{k+1}\le2(a\log(N/\varepsilon)+b)$,
for instance $N=Ms2^K$ with $K$ large, and not for every chain length; since the
block length $s2^{k+1}$ is even, no odd $N$ is covered. The source states the depth for every
$N$. Tensors that are not normal, for which the source combines the tree with
the preparation of a GHZ-type state, are treated without the tree
in~\cite{gap:mswc24_measurement_preparation_scope} and not here.

\section{Elimination}
The depth bound for every chain length of the circuit without measurements
uses unequal blocks: the last block has length $q'$ with $q\le q'<2q$
\cite{gap:mswc24_depth_upper_bound_divisible_length}. A complete binary tree of registers covers
$s2^{k+1}$ sites, so the last block needs a tree whose lowest level is unbalanced, or registers
of unequal lengths; each node is still a unitary on a bounded number of sites,
so the depth stays $O(k)$. This extends the trees of registers to blocks of
unequal lengths and is left open.

\bibliographystyle{alpha}
\bibliography{references}
\end{document}
